\section{The functional equation lemma}

In this section, we discuss the formalization of the functional equation
lemma. Throughout this section, all Lean declarations are made inside the
\lstinline|namespace FormalGroup|.  To simplify the presentation, we omit
the namespace keyword from the displayed code.
We begin by formalizing the data appearing in the statement.

Let $K$ be a commutative ring, let $R$ be a subring of $K$, and let $I$
be an ideal of $R$.  Let $p$, $q$, and $t$ be natural numbers, where $p$
is prime and $q$ is a power of $p$, say $q = p ^ t$, with $t \neq 0$.
Let $\sigma$ be a ring endomorphism of $K$, and let $s$ be a sequence of
elements of $K$.

In Lean, these assumptions are introduced as follows:
\begin{lstlisting}
variable {K : Type*} [CommRing K] {R : Subring K} {I : Ideal R}
  {p t q : ℕ} [hp : Fact (Nat.Prime p)] (ht : t ≠ 0) (hq : q = p ^ t)
  (σ : K →+* K) (s : ℕ → K) {g : PowerSeries R}
\end{lstlisting}

Recall that $f_g(X)$ is defined to be
the unique power series satisfying equation \eqref{equation: recurFunEq}.
We formalize $f_g(X)$ using the recursive formula of coefficient in
equation \eqref{eq: coeff_RecurFun} as

\begin{lstlisting}
def RecurFunAux (hg : g.constantCoeff = 0) : ℕ → K
  | 0 => 0
  | k + 1 =>
    g.coeff (k + 1) + ∑ j ∈ (Icc 1 (multiplicity q (k + 1))).attach,
      have : (k + 1) / (q ^ (j : ℕ)) < k + 1 := ...
      (s j) * σ^[j] (RecurFunAux hg ((k + 1) / (q ^ (j : ℕ))))

def RecurFun := PowerSeries.mk (RecurFunAux ht hq σ s hg)
\end{lstlisting}

% According to equation \eqref{eq: coeff_RecurFun}, we can inductively define a
% function \lstinline{RecurFunAux} and then obtain a power series
% using the API \lstinline{PowerSeries.mk}. We then prove that it satisfies
% the functional equation $\ref{equation: recurFunEq}$, namely
% \[f_g(X) = g(X) + \sum_{i=1}^{\infty} s_i \sigma_*^i f_g (X^{q ^ i}).\]
% \begin{lstlisting}
% def RecurFunAux (hg : g.constantCoeff = 0) : ℕ → K
%   | 0 => 0
%   | k + 1 =>
%     g.coeff (k + 1) + ∑ j ∈ (Icc 1 (multiplicity q (k + 1))).attach,
%       have : ((k + 1) / (q ^ (j : ℕ))) < k + 1 := ...
%       (s j) * σ^[j] (RecurFunAux hg ((k + 1) / (q ^ (j : ℕ))))

% def RecurFun := PowerSeries.mk (RecurFunAux ht hq σ s hg)
% \end{lstlisting}

% To inductively define \lstinline|RecurFunAux|, we prove that
% for all \lstinline|j ∈ Icc 1 (multiplicity q (k + 1))|,
% \lstinline|((k + 1) / (q ^ (j : ℕ))) < k + 1| in the motive of \lstinline|k + 1|. Therefore,
% we need to use \lstinline|Finset.attach|, which enables us to use the property
% \lstinline|j ∈ Icc 1 (multiplicity q (k + 1))|. To be precise,
% let \lstinline|s| be a finite set; \lstinline{Finset.attach s} takes the elements of \lstinline|s|
% and forms a new set of elements in the subtype \lstinline|{x // x ∈ s}|. With the property
% \lstinline|j ∈ Icc 1 (multiplicity q (k + 1))|, we can prove that
% \lstinline|((k + 1) / (q ^ (j : ℕ))) < k + 1|. Otherwise,
% \lstinline|(RecurFunAux hg ((k + 1) / (q ^ (j : ℕ))))| is not valid.

To define \lstinline|RecurFunAux| recursively, Lean must check that every
recursive call is made on a strictly smaller natural number.  In the
recursive step for \lstinline|k + 1|, the recursive call is made at
\lstinline|(k + 1) / (q ^ (j : ℕ))|, so we need a proof that
\lstinline|((k + 1) / (q ^ (j : ℕ))) < k + 1| for every index \lstinline|j|
appearing in the sum.
This is the reason for using \lstinline|Finset.attach|
\href{https://leanprover-community.github.io/mathlib4_docs/Mathlib/Data/Finset/Attach.html#Finset.attach}{\faExternalLink}
.
The summation is
taken over \lstinline|(Icc 1 (multiplicity q (k + 1))).attach|.
% \begin{lstlisting}
% (Icc 1 (multiplicity q (k + 1))).attach
% \end{lstlisting}
% rather than directly over \lstinline|Icc 1 (multiplicity q (k + 1))|.
An element of this attached finset is not just a natural number
\lstinline|j|, but a natural number together with a proof that it belongs
to the interval \lstinline|(Icc 1 (multiplicity q (k + 1)))|
% \begin{lstlisting}
% Icc 1 (multiplicity q (k + 1))
% \end{lstlisting}
In other words, from such a \lstinline|j| Lean can use the facts
\lstinline|1 ≤ j| and \lstinline|j ≤ multiplicity q (k + 1)|.

These bounds are needed to prove that division by \lstinline|q ^ j|
strictly decreases \lstinline|k + 1|.  Once Lean has the proof
\lstinline|(k + 1) / (q ^ (j : ℕ)) < k + 1|
% \begin{lstlisting}
% \end{lstlisting}
the recursive call
\lstinline|RecurFunAux hg ((k + 1) / (q ^ (j : ℕ)))|
% \begin{lstlisting}
% RecurFunAux hg ((k + 1) / (q ^ (j : ℕ)))
% \end{lstlisting}
is accepted as a structurally smaller call.  Without
\lstinline|Finset.attach|, the membership information about \lstinline|j|
would not be directly available inside the summation, and Lean would not
have enough information to justify the recursive definition.

Recall that functional equation \eqref{equation: recurFunEq} is
\[f_g(X) = g(X) + \sum_{i=1}^{\infty} s_i \sigma_*^i f_g (X^{q ^ i}).\]

And it is formalized as

\begin{lstlisting}
theorem FunEq_of_RecurFun [TopologicalSpace K] [T2Space K] (hs₀ : s 0 = 0) :
    let f := RecurFun ht hq σ s hg
    f = g.map R.subtype +
      ∑' (i : ℕ), s i • (f.expand (q ^ i) (q_pow_neZero hq)).map (σ^i) := ...
\end{lstlisting}

% The additional instance \lstinline|[TopologicalSpace K]| provides a topological structure on
% the ring of power series over $K$, and \lstinline|[T2Space K]| ensures that the
% infinite sum in the equation has a unique limit whenever it converges.
% Let $R$ be a subring of $K$. Then \lstinline|R.subtype|
% is a canonical ring homomorphism from $R$ to $K$. \lstinline|PowerSeries.expand p hp f| expands a power
% series \lstinline|f| by a nonzero factor $p$.

The assumptions \lstinline|[TopologicalSpace K]| and \lstinline|[T2Space K]|
are needed because the right-hand side contains the infinite sum
\lstinline|∑' (i : ℕ), ...|.  The topology on \lstinline|K| induces the
corresponding topology on \lstinline|PowerSeries K|, so that Lean can
interpret this expression as a limit of finite partial sums.  The
Hausdorff condition \lstinline|[T2Space K]| ensures that such limits are
unique, which is important when proving equality of power series defined
by an infinite summation.

In the formula, \lstinline|R.subtype| is the canonical ring homomorphism
from the subring \lstinline|R| to the ambient ring \lstinline|K|.  Thus
\lstinline|g.map R.subtype| regards the power series \lstinline|g| over
\lstinline|R| as a power series over \lstinline|K|.
The term \lstinline|(f.expand (q ^ i) (q_pow_neZero hq)).map (σ^i)|
corresponds to the expression $\sigma_*^i f_g (X^{q ^ i})$ in the
functional equation.  A discussion of \lstinline|MvPowerSeries.expand|
and its one-variable version for power series can be found in
Section~\ref{sec:mvpowerseries}.

Since the coefficient of $X$ in $f_g(X)$ is invertible in $K$.
Then we could the definition \lstinline|substInv| to define its substitution
inverse as

\begin{lstlisting}
def inv_RecurFun := (RecurFun ht hq σ s hg).substInv
\end{lstlisting}

% The theorem \lstinline|coeff_one_RecurFun| states that the first coefficient of $g$
% equals the first coefficient of $f_g(X)$. With the assumption
% \lstinline|(hg_unit : IsUnit (g.coeff 1))|,
% we obtain \lstinline|IsUnit ((RecurFun ht hq σ s hg).coeff 1)|. Therefore, there is a substitution inverse
% of $f_g(X)$, denoted by $f_g^{-1}(X)$.

% \begin{lstlisting}
% lemma coeff_one_RecurFun: (RecurFun ht hq σ s hg).coeff 1 = g.coeff 1 := ...

% lemma coeff_RecurFun_unit (hg_unit : IsUnit (g.coeff 1)) :
%     IsUnit ((RecurFun ht hq σ s hg).coeff 1) := ...

% def inv_RecurFun := PowerSeries.subst_inv _ (coeff_RecurFun_unit ht hq σ s hg hg_unit)
%   (constantCoeff_RecurFun ..)
% \end{lstlisting}

% Furthermore, let $F_g(X,Y) = f^{-1}_g(f_g{X} + f_g(Y))$ and $F_g(X,Y)$ is
% a commutative formal group law over $K$. And according to the functional equation lemma
% \ref{thm: fun_eq}, $F_g(X,Y)$ is a commutative formal group law over $R$.

The following definition constructs the bivariate power series
corresponding to $f_g^{-1}(f_g(X) + f_g(Y))$:
\begin{lstlisting}
def inv_add_RecurFun :=
    (inv_RecurFun ht hq σ s hg hg_unit).subst
      ((RecurFun ht hq σ s hg).toMvPowerSeries (0 : Fin 2) +
        (RecurFun ht hq σ s hg).toMvPowerSeries 1)
\end{lstlisting}

Here \lstinline|inv_RecurFun ht hq σ s hg hg_unit| denotes the
compositional inverse of the power series \lstinline|RecurFun ht hq σ s hg|.
The expression
\lstinline|(RecurFun ht hq σ s hg).toMvPowerSeries (0 : Fin 2)|
regards this one-variable power series as a bivariate power series in the
first variable, so it represents $f_g(X)$.  Similarly,
\lstinline|(RecurFun ht hq σ s hg).toMvPowerSeries 1| represents
$f_g(Y)$, where the second variable is used instead.

Therefore, the whole expression defines the bivariate power series
\[
  F_g(X,Y) = f_g^{-1}(f_g(X) + f_g(Y)).
\]

The first part of the functional equation lemma \ref{thm: fun_eq}
states that $F_g(X,Y)$ actually has coefficients in the subring $R$.
In Lean, this is
formalized by proving that every coefficient of
\lstinline|inv_add_RecurFun| belongs to \lstinline|R|:
\begin{lstlisting}
theorem coeff_inv_add_mem_Subring [UniformSpace K] [T2Space K] [DiscreteUniformity K]
    (hs₀ : s 0 = 0) :
    ∀ n, (inv_add_RecurFun ht hq σ s hg hg_unit).coeff n ∈ R := ...
\end{lstlisting}

After this theorem, the power series $F_g(X,Y)$ gives the desired
commutative formal group law over $R$.


\begin{lstlisting}
def inv_add_RecurFun :=
    (inv_RecurFun ht hq σ s hg hg_unit).subst
      ((RecurFun ht hq σ s hg).toMvPowerSeries (0 : Fin 2) +
        (RecurFun ht hq σ s hg).toMvPowerSeries 1)
\end{lstlisting}

The \lstinline|def inv_add_RecurFun| denotes the multivariate power series
$F_g(X,Y) = f_g^{-1}(f_g(X) + f_g(Y))$. Also,
\lstinline|(RecurFun ht hq σ s hg).toMvPowerSeries (0 : Fin 2)| represents the multivariate
power series $f_g(X)$, where $X$ is the first indeterminate variable in a bivariate power
series. Similarly, \lstinline|(RecurFun ht hq σ s hg).toMvPowerSeries 1| represents $f_g(Y)$.

Recall that the first part of the functional equation lemma $\ref{thm: fun_eq}$ says that
$F_g(X,Y)$ has coefficients in $R$. Moreover, it is a commutative formal group law by
the definition of $F_g(X,Y)$. The first part of the functional equation lemma is formalized as follows.

\begin{lstlisting}
theorem coeff_inv_add_mem_Subring [UniformSpace K] [T2Space K] [DiscreteUniformity K]
    (hs₀ : s 0 = 0) : ∀ n, (inv_add_RecurFun ht hq σ s hg hg_unit).coeff n ∈ R := ...
\end{lstlisting}

Next, we prove that $F_g(X,Y)$ is a formal group law over $K$. The associativity condition
follows from
\[F_g(F_g(X,Y), Z) = f_g^{-1} (f_g(X) + f_g(Y) + f_g(Z))=  F_g(X, F_g(Y,Z))\]

% We now explain how the functional equation lemma \ref{thm: fun_eq} is
% used to construct a formal group law over the subring $R$.

% First, the expression
% \[
%   F_g(X,Y) = f_g^{-1}(f_g(X) + f_g(Y))
% \]
% is naturally constructed as a power series over the larger ring $K$.
% Indeed, both \lstinline|RecurFun ht hq σ s hg| and its compositional
% inverse are power series over $K$.  Therefore, the first step is to define
% a formal group law over $K$.

% \begin{lstlisting}
% def invAdd_RecurFun_Aux : FormalGroup K where
%   toFun := inv_add_RecurFun ht hq σ s hg hg_unit
%   ...
% \end{lstlisting}

% Here \lstinline|invAdd_RecurFun_Aux| is the formal group law over $K$
% whose underlying power series is
% \lstinline|inv_add_RecurFun ht hq σ s hg hg_unit|, namely the power
% series representing $f_g^{-1}(f_g(X) + f_g(Y))$.

% The first part of the functional equation lemma \ref{thm: fun_eq} proves
% that all coefficients of this power series actually lie in the subring
% $R$:

% \begin{lstlisting}
% theorem coeff_inv_add_mem_Subring [UniformSpace K] [T2Space K] [DiscreteUniformity K]
%     (hs₀ : s 0 = 0) : ∀ n, (inv_add_RecurFun ht hq σ s hg hg_unit).coeff n ∈ R := ...
% \end{lstlisting}

Using this coefficient result, we can descend the formal group law from
$K$ to $R$.  This is done by applying \lstinline|FormalGroup.toSubring|
and the theorem \lstinline|coeff_inv_add_mem_Subring|.

\begin{lstlisting}
def invAdd_RecurFun [UniformSpace K] [T2Space K]
    [DiscreteUniformity K] (hs₀ : s 0 = 0) : FormalGroup R :=
  (FormalGroup.invAdd_RecurFun_Aux ht hq σ s hg hg_unit).toSubring _
    (coeff_inv_add_mem_Subring ...)
\end{lstlisting}

Moreover, this formal group law is commutative.  This follows from the
definition of $F_g(X,Y)$, since the expression
$f_g(X) + f_g(Y)$ is symmetric in $X$ and $Y$.

Thus the first part of the functional equation lemma \ref{thm: fun_eq}
gives a commutative formal group law over $R$, constructed from
the recursive power series $f_g$.

The second part of the functional equation lemma \ref{thm: fun_eq}
concerns the comparison between two recursively defined power series.
Suppose that $g$ and $g'$ are power series over $R$ with zero constant
coefficient.  The lemma says that the relevant composition involving the
inverse of one recursive series and the other recursive series again has
coefficients in $R$.

In Lean, this is formalized as follows:

\begin{lstlisting}
lemma coeff_inv_RecurFun_g'_mem_Subring [UniformSpace K] [T2Space K]
    [DiscreteUniformity K] (hs0 : s 0 = 0) :
    let f_g' := (RecurFun ht hq σ s hg')
    let G := (inv_RecurFun ht hq σ s hg hg_unit).subst f_g'
    ∀ n, PowerSeries.coeff n G ∈ R
\end{lstlisting}

Here \lstinline|f_g'| denotes the recursive power series associated to
\lstinline|g'|.  The power series \lstinline|G| is obtained by substituting
\lstinline|f_g'| into the inverse recursive power series.  The conclusion
states that every coefficient of \lstinline|G| lies in the subring
\lstinline|R|.

The third part of the functional equation lemma \ref{thm: fun_eq} gives a
stability result under substitution.  Let $h$
be a power series over $R$ with zero constant coefficient.  The goal is to show
that there exists a power series $\hat{h}$ over $R$ such that
\[
  f_g(h(X)) = f_{\hat{h}}(X).
\]

Using the functional equation \eqref{equation: recurFunEq}, this is equivalent to showing that the
power series
\[
  f_g(h(X)) -
  \sum_{n=1}^{\infty}
    s_n \cdot \sigma^n f_g(h(X^{q^n}))
\]
has coefficients in $R$.  This power series is then the desired
$\hat{h}$.

The corresponding Lean statement is:

\begin{lstlisting}
lemma coeff_inv_RecurFun_g'_mem_Subring' [UniformSpace K] [T2Space K]
    [DiscreteUniformity K] (hs₀ : s 0 = 0) {h : PowerSeries R}
    (h₀ : h.constantCoeff = 0):
    let f := (RecurFun ht hq σ s hg)
    let f₁ : PowerSeries K := f.subst (h.map R.subtype)
    ∀ n, (f₁ - ∑' i, (s i) • (f₁.expand (q ^ i) (q_pow_neZero hq)).map (σ ^ i)).coeff n ∈ R :=
\end{lstlisting}

In this statement, \lstinline|f₁| denotes the composition $f_g(h(X))$,
regarded as a power series over $K$.  The infinite sum is the correction
term appearing in the functional equation.  The conclusion says that the
coefficients of the resulting difference all lie in $R$.

The fourth part of the functional equation lemma \ref{thm: fun_eq}
proves a congruence statement.  It shows that congruence modulo a power of
the ideal $I$ is preserved after substituting into the recursive power
series $f_g$.

\begin{lstlisting}
theorem congr_equiv [UniformSpace K] [T2Space K] [DiscreteUniformity K]
    (hs₀ : s 0 = 0) {α : PowerSeries R} {β : PowerSeries K}
    (hα : α.constantCoeff = 0) (hβ : β.constantCoeff = 0)
    {r : ℕ} (hr : 1 ≤ r) :
    let f := RecurFun ht hq σ s hg
    (∀ n, α.coeff n - β.coeff n ∈ R.subtype '' ↑(I ^ r)) ↔
      (∀ n, PowerSeries.coeff n (f.subst (α.map R.subtype)) -
        PowerSeries.coeff n (f.subst β) ∈ R.subtype '' ↑(I ^ r)) :=
\end{lstlisting}

There is a small technical point in this statement.  The power series
\lstinline|α| is defined over $R$, while \lstinline|β| is defined over
$K$.  Therefore, after comparing their coefficients, the difference
\lstinline|α.coeff n - β.coeff n| is regarded as an element of $K$.
For this reason, the congruence condition is expressed as membership in a
subset of $K$, rather than directly as membership in the ideal
\lstinline|I ^ r| of $R$.

The subset \lstinline|R.subtype '' ↑(I ^ r)|
% \begin{lstlisting}
% R.subtype '' ↑(I ^ r)
% \end{lstlisting}
is the image of the ideal \lstinline|I ^ r| under the natural inclusion
\lstinline|R.subtype : R →+* K|.  Thus it consists of those elements of
$K$ which come from elements of \lstinline|I ^ r|.  The symbol
\lstinline|↑| is needed because Lean coerces the ideal
\lstinline|I ^ r : Ideal R| to its underlying set \lstinline|I ^ r : Set R|
before taking its image under
\lstinline|R.subtype|.

Hence the theorem says that two power series are coefficientwise
congruent modulo \lstinline|I ^ r| before applying \lstinline|f_g| if and
only if their images under substitution into \lstinline|f_g| are still
coefficientwise congruent modulo \lstinline|I ^ r|.

The full proofs of these statements can be found in the GitHub
repository \href{https://github.com/WenrongZou/FormalGroupLaws/blob/main/FormalGroupLaws/IntegralityLem.lean}{\faExternalLink}.

% We now explain how the functional equation lemma is used to construct a
% formal group law over the subring $R$.

% First, the expression
% \[
%   F_g(X,Y) = f_g^{-1}(f_g(X) + f_g(Y))
% \]
% is naturally constructed as a power series over the larger ring $K$.
% Indeed, both \lstinline|RecurFun ht hq σ s hg| and its compositional
% inverse are power series over $K$.  Therefore, the first step is to define
% a formal group law over $K$.

% Then we use the API \lstinline|FormalGroup.toSubring| and the theorem
% \lstinline|coeff_inv_add_mem_Subring| to obtain a formal group law
% \lstinline|FormalGroup.InvAdd_RecurFun| over $R$.

% \begin{lstlisting}
% def invAdd_RecurFun_Aux : FormalGroup K where
%   toFun := inv_add_RecurFun ht hq σ s hg hg_unit
%   ...

% def invAdd_RecurFun [UniformSpace K] [T2Space K]
%     [DiscreteUniformity K] (hs₀ : s 0 = 0) : FormalGroup R :=
%   (FormalGroup.invAdd_RecurFun_Aux ht hq σ s hg hg_unit).toSubring _
%     (coeff_inv_add_mem_Subring ...)

% instance : (InvAdd_RecurFun ..).IsComm := sorry
% \end{lstlisting}

% Recall that the second part of the functional equation lemma $\ref{thm: fun_eq}$ states that,
% for any power series $g, g'$ over $R$ with the assumption that the constant coefficients of $g$ and $g'$
% are zero, $f_{g'}^{-1} (f_g(X))$ has coefficients in $R$. The second part of the functional equation lemma
% is formalized as follows. In the lemma below, for simplicity, we denote $f_{g'}(X)$ by
% \lstinline|f_g'| and $f_{g'}^{-1}(f_g(X))$ by \lstinline|G|.

% \begin{lstlisting}
% lemma coeff_inv_RecurFun_g'_mem_Subring [UniformSpace K] [T2Space K]
%     [DiscreteUniformity K] (hs0 : s 0 = 0) :
%     let f_g' := (RecurFun ht hq σ s hg')
%     let G := (inv_RecurFun ht hq σ s hg hg_unit).subst f_g'
%     ∀ n, PowerSeries.coeff n G ∈ R
% \end{lstlisting}

% The third part of the functional equation lemma \ref{thm: fun_eq} states that there is a power series $\hat{h}$ over $R$
% such that $f_g(h(X)) = f_{\hat{h}}(X)$. According to the functional equation
% $\ref{equation: recurFunEq}$, this is equivalent to showing that the power series
% $f_g(h(X)) - \sum_{n= 1}^{\infty} s_n \cdot \sigma^n f_g(h(X^{q^n}))$ is a power series over $R$.
% This is exactly the desired power series $\hat{h}$.

% \begin{lstlisting}
% lemma coeff_inv_RecurFun_g'_mem_Subring' [UniformSpace K] [T2Space K]
%     [DiscreteUniformity K] (hs₀ : s 0 = 0) {h : PowerSeries R}
%     (h₀ : h.constantCoeff = 0):
%     let f := (RecurFun ht hq σ s hg)
%     let f₁ : PowerSeries K := f.subst (h.map R.subtype)
%     ∀ n, (f₁ - ∑' i, (s i) • (f₁.expand (q ^ i) (q_pow_neZero hq)).map (σ ^ i)).coeff n ∈ R :=
% \end{lstlisting}

% \begin{lstlisting}
% theorem congr_equiv [UniformSpace K] [T2Space K] [DiscreteUniformity K]
%     (hs₀ : s 0 = 0) {α : PowerSeries R} {β : PowerSeries K}
%     (hα : α.constantCoeff = 0) (hβ : β.constantCoeff = 0)
%     {r : ℕ} (hr : 1 ≤ r) :
%     let f := RecurFun ht hq σ s hg
%     (∀ n, α.coeff n - β.coeff n ∈ R.subtype '' ↑(I ^ r)) ↔
%       (∀ n, PowerSeries.coeff n (f.subst (α.map R.subtype)) -
%         PowerSeries.coeff n (f.subst β) ∈ R.subtype '' ↑(I ^ r)) :=
% \end{lstlisting}

% I briefly explain some details in part four of the functional equation lemma \ref{thm: fun_eq}.
% To begin with, \lstinline|α.coeff n - β.coeff n| is an element of $K$, since $\alpha$ is
% a power series over $R$ and $\beta$ is a power series over $K$. Therefore, we need to formalize it
% as an element of a subset of $K$ rather than as an element of an ideal of $R$.
% Secondly, \lstinline|R.subtype '' ↑(I ^ r)| is a subset of $K$ that is the image of $I^r$
% under the natural ring homomorphism $R \to K$. Moreover, the symbol \lstinline|↑| is necessary
% since there is a coercion from \lstinline|Ideal R| to \lstinline|Set R|.

% The proof of the above theorem can be found in the GitHub repository.